% !TeX program = lualatex
% !TeX encoding = utf8
% !TeX spellcheck = uk_UA
% !BIB program = biber

\documentclass[%
biblatex,
%marginversioninfo
]{ConspectBook}

%========================================================================================================
%
%									      Вихідні дані докумета
%
%========================================================================================================

\title{Класична електродинаміка}
\def\subtitle{Мікроскопічна теорія}


%======================================== Бібліографія ===================================================
%
\addbibresource{d:/Projects/LaTeX/MyPackages/bib/MyBase.bib}
%\addbibresource{Kvant.bib}
%
%========================================================================================================

%========================================================================================================
%
% !TeX root =./FTProblems.tex

%========================================================================================================
%
%									      Палітурка
%
%========================================================================================================

\newcommand{\CoverPage}{
	\begin{alwayssingle}
		\begin{center}
			\begin{flushright}\bfseries\sffamily
				\MakeUppercase{Міністерство освіти і науки України}\\
				КПІ ім. Ігоря Сікорського\\
			\end{flushright}
			\begin{tcolorbox}[titlepagestyle,
					toprule=0.10cm,
					bottomrule=0.10cm,
					overlay={%
						\node (picture) at ([xshift=4cm]frame.west) {\includegraphics[width=3.5cm]{logo_IPT}};
					}
			]%
			\begin{flushright}
				\large\bfseries\color{white}Фізико-технічний інститут
			\end{flushright}
			\end{tcolorbox}
			\vspace*{-2em}
		 	\begin{Large}\color{themecolordark!90!black}
			\begin{alignat*}{3}
				\vect{\nabla}&\times&&\Efield &&= -\dfrac{1}{c}\dfrac{\partial\Bfield}{\partial t} \\
				\vect{\nabla}&\,\cdot&&\Bfield &&= 0 \\
				\vect{\nabla}&\,\cdot&&\Dfield &&= 4\pi\rho \\
				\vect{\nabla}&\times&&\Hfield &&= \dfrac{4\pi}{c} \vect{j}+\dfrac{1}{c}\dfrac{\partial\Dfield}{\partial t}
			\end{alignat*}
			\end{Large}
			\vspace*{2em}
			\begin{tcolorbox}[
				titlepagestyle,
				toprule=0.15cm,
				bottomrule=0.15cm,
				top=1.3cm,
				bottom=0.7cm,
				overlay={%
				\node[%
							fill=white,
							rounded corners = 15pt,
							draw=themecolorlight,
							line width=0.15cm,
							inner sep=0pt,
							text width=17cm,
							minimum height=2cm,
							align=center,
							%anchor=east,
							font=\sffamily\bfseries\large
						] (title) at (frame.north) {В.~І. Жданов, %
													С.~М. Пономаренко,
													В.~Б. Долгошей%
					};
				}
			]
			\centering
			\Huge\sffamily\bfseries\textcolor{white}{\realtitle}\\
			\huge\sffamily\bfseries\textcolor{white}{\subtitle}
			\end{tcolorbox}
			\vfill
		 	\begin{Large}\color{themecolordark!90!black}
			\begin{gather*}
				\frac{\partial F^{\mu\nu}}{\partial x^\nu} =                                                                                -\frac{4\pi}{c}j^\mu \\
				\frac{\partial F_{\mu\nu}}{\partial x^\alpha} + \frac{\partial F_{\mu\alpha}}{\partial x^\nu} +\frac{\partial F_{\alpha\mu}}{\partial x^\nu}  = 0.
			\end{gather*}
			\end{Large}
			\vfill
			\begin{tcolorbox}[titlepagestyle,
					toprule=0.10cm,
					bottomrule=0.10cm]
				\begin{center}\color{white}\bfseries\normalsize
					\MakeUppercase{Київ~\the\year} \\
%					КПІ ім. Ігоря Сікорського \\
%					\the\year
				\end{center}
			\end{tcolorbox}
		\end{center}
		\clearpage
	\end{alwayssingle}
\setcounter{page}{1}
}


%========================================================================================================
%
%									      Титульна сторінка
%
%========================================================================================================

\renewcommand\maketitle{
	\begin{alwayssingle}
		\begin{center}
				\MakeUppercase{Міністерство освіти і науки України}

				\bigskip
				\MakeUppercase{Національний технічний університет України}\\
				<<КИЇВСЬКИЙ ПОЛІТЕХНІЧНИЙ ІНСТИТУТ \\ імені ІГОРЯ СІКОРСЬКОГО>>
				\vspace*{100pt}

				{\large В.~І.~Жданов,
						С.~М.~Пономаренко,
						В.~Б.~Долгошей
						}
				\vspace*{50pt}

				{\Huge\sffamily\bfseries\realtitle}\\[1em]
				{\huge\sffamily\bfseries\subtitle}

			\vspace*{50pt}
			\begin{center}\itshape
			Рекомендовано Методичною радою КПІ ім. Ігоря Сікорського як навчальний посібник для здобувачів ступеня бакалавра за спеціальностями 105 <<Прикладна фізика та наноматеріали>>, 113 <<Прикладна математика>>
			\end{center}

			\vfill
			\begin{center}
				\MakeUppercase{Київ} \\
				КПІ ім. Ігоря Сікорського \\
				\the\year
			\end{center}
		\end{center}
		\clearpage
	\end{alwayssingle}
}


%========================================================================================================
%
%									      Друга сторінка
%
%========================================================================================================
\newcommand\makeinfopage{
	\begin{alwayssingle}
		\noindent%
		\begin{minipage}[t]{0.5\textwidth}
				\begin{flushleft}
					УДК  537\\
					Ж 42
				\end{flushleft}
		\end{minipage}


		\bigskip\noindent%
        \begin{minipage}[t]{0.2\linewidth}
            	\begin{flushleft}
                    Рецензенти:
                \end{flushleft}
        \end{minipage}\hfill
        \begin{minipage}[t]{0.78\linewidth}
                \href{http://www.nas.gov.ua/UA/PersonalSite/Pages/default.aspx?PersonID=0000006576}{В.~О.~Кочелап}, д.ф.-м.н., професор,
                член-кореспондент НАН України, завідувач відділу теоретичної фізики Інституту фізики напівпровідників ім. В.~Є. Лашкарьова
%                \\[1ex]
%                \href{http://apd.ipt.kpi.ua/blog/author/19}{Я.~Д.~Кривенко-Еметов}, к.ф.-м.н., доцент кафедри прикладної фізики, Фізико-технічного
%інституту КПІ ім. Ігоря Сікорського
        \end{minipage}

		\bigskip\noindent%
        \begin{minipage}[t]{0.2\linewidth}
            	\begin{flushleft}
                    Відповідальний редактор:
                \end{flushleft}
        \end{minipage}\hfill
        \begin{minipage}[t]{0.78\linewidth}
                \href{http://ipt.kpi.ua/litvinova}{Т.~В.~Литвинова}, к.т.н., доцент, перший заступник директора Фізико-технічного інституту КПІ ім. Ігоря Сікорського
        \end{minipage}

		\begin{center}\itshape\small
				Гриф надано Методичною радою КПІ ім. Ігоря Сікорського (протокол №~10/2020~від 18.06.2020~р.) за поданням Вченої ради Фізико-технічного інституту (протокол №~5/2020 від 25.05.2020 р.)
		\end{center}
		\begin{center}
			\ifelectronic Електронне мережне навчальне видання \fi
			%\par {Версія від~\href{http://www.istpravda.com.ua/dates}{\today}} \par\else \par  \fi
		\end{center}
		\begin{center}
			\href{http://phes.ipt.kpi.ua/zhdanov-valerij-ivanovych}{\itshape Жданов Валерій Іванович}, д.ф.-м.н., професор \\
			\href{http://phes.ipt.kpi.ua/ponomarenko-sergij-mikolajovich}{\itshape Пономаренко Сергій Миколайович}, к.ф.-м.н., доцент \\
			\href{http://phes.ipt.kpi.ua/dolgoshej-volodymyr-borysovych}{\itshape Долгошей Володимир Борисович}, к.ф.-м.н., доцент
		\end{center}
%			\vspace*{1em}%
		\begin{center}\bfseries
			\LARGE\sffamily\realtitle \\
			\Large\sffamily\subtitle
		\end{center}
        \noindent%
        \begin{minipage}[t]{\textwidth}\small
                \realtitle: \subtitle\ [Електронний ресурс] : навч. посіб. для студ. спеціальностей
                105 <<Прикладна фізика та наноматеріали>> та 113 <<Прикладна математика>> /  В.~І. Жданов, С.~М. Пономаренко, В.~Б. Долгошей ; КПІ ім. Ігоря Сікорського.~--- Електронні текстові дані
            (1 файл: 570~кБ). – Київ : КПІ ім. Ігоря Сікорського, \the\year. --- \the\numexpr\getpagerefnumber{LastPages}-1\relax~с.
        \end{minipage}

%		\noindent%
%		\begin{flushleft}
%			\begin{tabular}{lp{0.9\textwidth}}
%				     & В.~І. Жданов, С.~М. Пономаренко, В.~Б. Долгошей                                                                                                                                                                        \\
%				Ж 42 & \hspace*{3ex} \realtitle : \subtitle{} [Електронний ресурс] : навчальний посібник / %
%				В.~І. Жданов, %
%				С.~М. Пономаренко%,
%				В.~Б. Долгошей %
%				--- К.:~КПІ ім. Ігоря Сікорського, \the\year. --~\the\numexpr\getpagerefnumber{LastPage}-1\relax~с. -- Бібліогр.: с.~\pageref{BibPage}\relax. \ifelectronic\relax\else-- 80~прим.\fi
%			\end{tabular}
%		\end{flushleft}
		\vfill

		Наведено \TotalValue{totalproblems} задач з курсу класичної електродинаміки. До більшості задач подано відповіді. Задачі різноманітні як за змістом, так і за складністю. Поряд із задачами, які ілюструють основні поняття і закони електродинаміки, до збірника включена певна кількість більш складних задач, які допомагають суттєвішому вивченню предмета.

		Для студентів фізико-технічного інституту КПІ ім. Ігоря Сікорського, які навчаються за спеціальностями 105~<<Прикладна фізика та наноматеріали>> та 113 <<Прикладна математика>>.

		\vfill

%		\begin{flushleft}\small
%			Ілюстративний матеріал підручника підготовлений за допомогою пакету \href{http://pgf.sourceforge.net}{TikZ/Pgf}. Верстка тексту проведена в видавничій системі \LaTeXe{} (компілятор Lua\LaTeX) на базі системи комп'ютерної верстки \TeX{} (Збірка  \href{https://www.tug.org/texlive/}{\TeX Live~\the\year}) з використанням оболонки \href{https://www.texstudio.org}{\TeX Studio}.
%		\end{flushleft}
	\hfill
	\begin{minipage}[t]{0.65\linewidth}\small
        \textcopyright{} В.~І. Жданов, С.~М. Пономаренко, В.~Б. Долгошей, \the\year\,р. \\
        \textcopyright{}  КПІ ім. Ігоря Сікорського (ФТІ), \the\year~р.
    \end{minipage}
		\newpage%
	\end{alwayssingle}
}

%%%%%%%%%%%%%%%%%%%%%%%%%%%%%%%%%%%%%%%%%%%%%%%%%%%%%%%%%%%%%%%%%%%%%%%%%%%%
%%                                                                        %%
%%                              Last Page                                 %%
%%                                                                        %%
%%%%%%%%%%%%%%%%%%%%%%%%%%%%%%%%%%%%%%%%%%%%%%%%%%%%%%%%%%%%%%%%%%%%%%%%%%%%
\newcommand{\makelastpage}{%
\clearpage%
\thispagestyle{empty}%
\vspace*{0.4\textheight}
\begin{center}
		\textbf{Жданов} Валерій Іванович \\
	    \textbf{Пономаренко} Сергій Миколайович \\
	    \textbf{Долгошей} Володимир Борисович
\end{center}

\begin{center}\bfseries
    \Large\sffamily\realtitle \\
    \large\sffamily\subtitle
\end{center}

\begin{center}\itshape
    Комп'ютерне верстання в системі \LaTeXe{} С.\ М. Пономаренко
\end{center}

\vspace*{1em}
\begin{center}\small
Національний технічний університет України \\
<<Київський політехнічний інститут імені Ігоря Сікорського>> \\
Свідоцтво про державну реєстрацію: серія ДК № 5354 від 25.05.2017 р.\\
просп. Перемоги, 37, м. Київ, 03056
\end{center}

%\vfill
%
%\begin{center}
%    Видавництво <<Політехніка>> \kpishort \\
%    вул. Політехнічна, 14, корп. 15 \\
%    м. Київ, 03056 \\
%    Тел. (044) 204-81-78
%\end{center}
}


%					      Титульна сторінка
%
%========================================================================================================

\begin{document}
\pagestyle{empty}
%\ifelectronic\CoverPage\setcounter{page}{0}\fi
\CoverPage
\pagenumbering{arabic}
\maketitle
%\makeinfopage
%========================================================================================================
%
\clearpage\pagestyle{plain}
\tableofcontents%				               Зміст
%
%========================================================================================================

%========================================================================================================
%
%% !TeX program = lualatex
% !TeX encoding = utf8
% !TeX spellcheck = uk_UA
% !TeX root =../EMProblems.tex

\introtrue
\chapter*{Передмова}

Вивчення класичної електродинаміки забезпечується дисциплінами <<Тео\-рія поля>> та <<Електродинаміка суцільних середовищ>>, які входять до циклу базової та професійної підготовки бакалаврів за спеціальностями 105 <<Прикладна фізика та наноматеріали>> і 113 <<Прикладна математика>> й вивчаються студентами Фізико-технічного інституту КПІ ім. Ігоря Сікорського на третьому курсі. На вивчення дисциплін відведено два семестри, практичні заняття проводяться один раз на тиждень. Велика частина програмного матеріалу пов'язана з умінням розв'язування конкретних задач. Вироблення умінь, навичок і методів розв'язку величезної кількості задач, звичайно, не може бути реалізоване тільки за рахунок годин, відведених на практичні заняття, і передбачає значну самостійну роботу студента.

Назви розділів та підрозділів посібника відповідають робочим програмам дисциплін. До кожного з них подано короткий теоретичний матеріал, який містить основні формули, необхідні для розв'язування задач. В кінці збірника міститься довідковий матеріал та перелік використаної літератури. 


Задачі даного збірника щорічно пропонуються студентам при складанні письмового екзамену, який передбачено навчальним планом. Зрозуміло, що багато задач було взято з джерел, які перелічені в списку літератури, оскільки методика викладання дисциплін протягом  століть вже є переважно завершеною і багато вже винайдених іншими вченими задач є ключовими для розуміння предмета. Однак в збірнику є й велика частина авторських задач, ідеї яких сформувались на основі власного досвіду викладання предмету.

Як правило, в тексті використовуються такі позначення координат: $x$,$y$,$z$ для декартових координат;  $r$, $\theta$, $\phi$ для сферичних та  $r$, $\phi$, $z$ для циліндричних (в останніх позначеннях використано $r=\sqrt{x^2+y^2}$,  щоб уникнути плутанини з густиною заряду).
 
Дана версія посібника є електронним виданням, тому для зручності користування ним передбачена система навігації у вигляді гіперпосилань та бічної панелі змісту.

%\vspace*{4em}
%
%\begin{flushright}\Annabelle
%	С.~М.~Пономаренко
%\end{flushright}
\introfalse





%                          Вступ та передмова
%
%========================================================================================================

%========================================================================================================
%
%									      Вставка файлів розділів
%
%========================================================================================================


%% !TeX program = lualatex
% !TeX encoding = utf8
% !TeX spellcheck = uk_UA
% !TeX root =../FTProblems.tex

\newpage
\chapter{Система \mbox{мікроскопічних} рівнянь Максвелла}

%\epigraph{\Annabelle  Неможливо позбутися відчуття, що ці математичні формули існують незалежно від нас і володіють власним розумом, що вони мудріші за нас, мудріше навіть тих, хто їх відкрив, і що ми дістаємо з них більше, ніж спочатку було закладено\ldots}{Heinrich Hertz}
%\epigraph{\Annabelle Уж не Боги ли начертали эти знаки?}{Л.~Больцман}

%\renewcommand{\theequation}{\thepart.\arabic{equation}}
\section{Рівняння Максвелла}

Система рівнянь Максвелла сформульована в 60-х роках XIX століття на основі узагальнення емпіричних законів і розвитку ідей про електромагнітні явища. Система рівнянь є основою теорії електромагнітного поля і дозволяє розв'язувати задачі, пов'язані з відшуканням електричних і магнітних полів, що створюються заданим розподілом електричних зарядів і струмів.

Інтегральна форма мікроскопічних рівнянь:
\begin{align}
	\oiint\limits_{\partial\Omega} \Efield\cdot d\vect{S} & = 4\pi\iiint\limits_{\Omega}\rho dV   \label{Int
	I},                                                                                                         \\
	\oiint\limits_{\partial\Omega} \Bfield\cdot d\vect{S} & = 0   \label{Int
	II},                                                                                                                                   \\
	\oint\limits_{\partial S} \Efield\cdot d\vect{r}  & = - \frac1c \iint\limits_S \frac{\partial\Bfield}{\partial t}\cdot d\vect{S}  \label{Int
	III},                                                          \\
	\oint\limits_{\partial S} \Bfield\cdot d\vect{r}  & =\dfrac{4\pi}{c} \iint\limits_S \vect{j}\cdot d\vect{S} +\frac{1}{c} \iint\limits_S
	\frac{\partial\Efield}{\partial t}\cdot d\vect{S}  \label{Int IV},
\end{align}
де $\Omega$ --- довільний нерухомий об'єм, $\partial\Omega$ --- його межа; $S$ --- довільна нерухома орієнтована поверхня, $\partial S$ --- замкнений
контур, що її обмежує.

Диференціальна форма  мікроскопічних рівнянь:
\begin{flalign}
	\divg\Efield &= 4\pi\rho \label{Diff I},\\[0.8em]
	\divg\Bfield &= 0 \label{Diff II},\\
	\rot\Efield &= -\dfrac{1}{c}\dfrac{\partial\Bfield}{\partial t} \label{Diff III},\\
	\rot\Bfield &= \dfrac{4\pi}{c} \vect{j}+\dfrac{1}{c}\dfrac{\partial\Efield}{\partial t} \label{Diff IV}.
\end{flalign}


\noindent%
$ \Efield $~-- вектор напруженості електричного поля,\\
$ \Bfield $~-- вектор індукції магнітного поля.


\section{Закони збереження}

\begin{itemize}
	\item Електричного заряду
	      \begin{equation}
		      \frac{\partial \rho}{\partial t} + \divg\vect{j} = 0.
	      \end{equation}
	\item Збереження енергії (є наслідком рівнянь Максвелла)
	      \begin{equation}
		      \frac{\partial w}{\partial t} + \divg\vect{S} = -  \vect{j}\cdot\Efield,
	      \end{equation}
\end{itemize}
де $ w $~-- густина енергії електромагнітного поля:
\begin{equation}
	w = \frac{1}{8\pi} (\Efield^2 + \Bfield^2),
\end{equation}
\noindent%
$ \vect{S} $~-- вектор Пойнтінга (вектор густини потоку енергії):
\begin{equation}
	\vect{S} = \frac{c}{4\pi} \Efield\times\Bfield.
\end{equation}
Величина
\begin{equation}
	p =\vect{j}\cdot\Efield ,
\end{equation}
є потужністю, що виділяється в одиниці об'єму речовини.

\section{Граничні умови для границі розділу двох середовищ}

\begin{align}
	\left[ \vect{n} \times (\Efield_{1} - \Efield_{2}) \right] & = 0, \label{EnBc}                        \\
	\left( \Efield_{1} - \Efield_{2}\right) \cdot \vect{n}     & = 4\pi\sigma, \label{EtauBc}             \\
	\left( \Bfield_{1} - \Bfield_{2}\right) \cdot \vect{n}     & = 0,    \label{BnBc}                     \\
	\left[ \vect{n} \times (\Bfield_{1} - \Bfield_{2}) \right] & = \frac{4\pi}{c}\vect{i} \label{BtauBc},
\end{align}

\noindent%
де $\sigma$~-- поверхнева густина вільних зарядів на границі розділу,\\
$\vect{i}$~-- поверхневий струм провідності, який протікає по границі розділу, \\
$\vect{n}$~--- вектор нормалі з середовища $2$  в середовище $1$.

\section{Потенціали електромагнітного поля}

\begin{align}
	\Efield & = -\frac{1}{c}\frac{\partial \vect{A}}{\partial t} - \vect{\nabla}\phi, \label{E=phi+A} \\
	\Bfield & = \rot\vect{A},
\end{align}
де $\phi$~-- скалярний потенціал, $ \vect{A} $~-- вектор-потенціал.

Потенціали електромагнітного поля не визначені однозначно. Якщо замість потенціалів $\vect{A}$ та $\phi$ вибрати інші $\vect{A}'$ та $\phi'$, які пов'язані з вихідними перетвореннями:
\begin{align}
	\vect{A}' & = \vect{A} +\vect{\nabla}f,                     \\
	\phi'     & = \phi - \frac1c \frac{\partial f}{\partial t},
\end{align}
де  $f(x,y,z,t)$~-- довільна функція, то поля $\Efield$ та $\Bfield$ залишаються незмінними.

Умова
\begin{equation}\label{Lorentz_gauge}
	\vect\nabla\cdot\vect{A} = - \frac{1}{c}\frac{\partial\phi}{\partial t},
\end{equation}
називається калібруванням Лоренца і зазвичай використовується для динамічних задач.
Використовуючи калібрування Лоренца, рівняння для скалярного і векторного потенціалів можна записати у симетричному вигляді:
\begin{align}\label{Maxwell_equations_via_potantials}
	\Delta \vect{\phi} - \frac{1}{c^2} \frac{\partial^2 \phi}{\partial t^2} = -4\pi\rho, \\
	\Delta \vect{A} - \frac{1}{c^2} \frac{\partial^2 \vect{A}}{\partial t^2} = - \frac{4\pi}{c}\vect{j}.
\end{align}

У випадку ізольованої системи зарядів та струмів розв'язками цих рівнянь є \emph{запізнюючі потенціали}:
\begin{equation}\label{Retarded_potentials}
	\phi(\vect{r} , t) = \iiint\limits_{\Omega}\frac{\rho\left( \vect{r}', t - \frac{|\vect{r} - \vect{r}'|}{c}\right) }{|\vect{r} - \vect{r}'|}dV',
	\quad
	\vect{A}(\vect{r} , t) = \frac1c\iiint\limits_{\Omega}\frac{\vect{j}\left( \vect{r}', t - \frac{|\vect{r} - \vect{r}'|}{c}\right) }{|\vect{r} -
	\vect{r}'|}dV'.
\end{equation}

%\renewcommand{\theequation}{\thechapter.\arabic{equation}}

\multiinclude{
    Basics,
    Solutions,
    FreeField,
    Radiation
}[]


%========================================================================================================
%
%									       Вставка відповідей
%
%========================================================================================================
%\bookmarksetup{startatroot}
%\answers
%\multiinclude{\ChaptersOne}[-Answers]
%\multiinclude{\ChaptersTwo}[-Answers]
%========================================================================================================
%
%									             Додатки
%
%========================================================================================================
\appendix
% !TeX program = lualatex
% !TeX encoding = utf8
% !TeX spellcheck = uk_UA
% !TeX root =../FTProblems.tex

\section{Основні формули векторного аналізу}\label{Vanaliz}

\subsection{Тривимірний символ Леві-Чівіти}\label{Levi-Chiv}

Тривимірний символ Леві-Чівіти визначений співвідношеннями
\begin{equation}\label{eq:Levi-Chiv}
\epsilon_{ijk} =  - \epsilon_{jik} =  - \epsilon_{ikj}.
\end{equation}

Деякі співвідношення з цим символом:
\begin{align}
    \epsilon_{ijk}\epsilon_{kpq} &= \delta_{ip}\delta_{jq} - \delta_{iq}\delta_{jp}, \\
    \epsilon_{iqk}\epsilon_{pqk} &= 2\delta _{ip}.
\end{align}

Ротор та векторний добуток у декартових координатах $\{x,y,z\} = \{x_1, x_2, x_3\}$:
\begin{align}
\left[\rot\vect{A} \right]_i &= \epsilon_{ijk}{\partial_j}{A_k}, \\
\left[ \vect{A} \times \vect{B} \right]_i &= \epsilon_{ijk}{A_j}{B_k},
\end{align}
всі індекси пробігають значення $1$, $2$, $3$.

\subsection{Диференціальні операції в різних системах координат}

\subsubsection{Декартова система координат}

\begin{align}\label{cartesian}
	\mathrm{grad}\,\psi \equiv \vect{\nabla}\psi & = \frac{\partial \psi}{\partial x} \vect{e}_x + \frac{\partial \psi}{\partial y} \vect{e}_y +
	\frac{\partial \psi}{\partial z} \vect{e}_z \\
	\mathrm{div}\,\left(\mathrm{grad}\,\psi\right) \equiv \Laplasian\psi    & = \frac{\partial^2 \psi}{\partial x^2} + \frac{\partial^2 \psi}{\partial y^2} + \frac{\partial^2 \psi }{\partial z^2}               \\
	\mathrm{div}\,\vect{A} \equiv \divg\vect{A}     & = \frac{\partial A_x}{\partial x}  + \frac{\partial A_y}{\partial y} + \frac{\partial A_z}{\partial z}                              \\
	\mathrm{rot}\,\vect{A} \equiv  \rot\vect{A}      & = \left( \frac{\partial A_z}{\partial y}  - \frac{\partial A_y}{\partial z}\right)  \vect{e}_x +
	\left( \frac{\partial A_x}{\partial z}  - \frac{\partial A_z}{\partial x}\right)  \vect{e}_y +
	\left( \frac{\partial A_y}{\partial x}  - \frac{\partial A_x}{\partial y}\right)  \vect{e}_z
\end{align}

\subsubsection{Циліндрична система координат}

\begin{align}\label{cylindric}
		\vect{\nabla}\psi & = \frac{\partial \psi}{\partial \rho} \vect{e}_{\rho} + \frac{1}{\rho}\frac{\partial \psi}{\partial \phi} \vect{e}_{\phi} + \frac{\partial \psi}{\partial z} \vect{e}_{z}                                                                                       \\
	\Laplasian\psi    & = \frac{1}{r} \frac{\partial }{\partial r} \left( r \frac{\partial \psi}{\partial r} \right) + \frac{1}{r^2} \frac{\partial^2 \psi}{\partial \phi^2} + \frac{\partial^2 \psi}{\partial z^2}                                                                   \\
	\divg\vect{A}     & = \frac{1}{\rho}\frac{\partial \left(\rho A_{\rho }\right)}{\partial \rho }+\frac{1}{\rho }\frac{\partial A_{\phi } }{\partial \phi }+\frac{\partial A_{z}}{\partial z}                                                                                    \\
	\rot\vect{A}      & =\left({\frac {1}{\rho }}{\frac {\partial A_{z}}{\partial \phi }}-{\frac {\partial A_{\phi }}{\partial z}}\right)\vect{e}_{\rho}+\left({\frac {\partial A_{\rho }}{\partial z}}-{\frac {\partial A_{z}}{\partial \rho }}\right)\vect{e}_{\phi} + \nonumber \\
	                  & +
	{\frac {1}{\rho }}\left({\frac {\partial (\rho A_{\phi })}{\partial \rho }}-{\frac {\partial A_{\rho }}{\partial \phi }}\right)\vect{e}_{z}
\end{align}
Орти циліндричної системи координат зв'язані з декартовими ортами як:
\begin{align}\label{}
    \vect{e}_r &= \cos\phi\, \vect{e}_x + \sin\phi\, \vect{e}_y, \\
    \vect{e}_{\phi} &= -\sin\phi\, \vect{e}_x + \cos\phi\, \vect{e}_y, \\
    \vect{e}_{z} &= \vect{e}_z.
\end{align}

\subsubsection{Сферична система координат}

\begin{align}\label{spheric}
    \vect{\nabla}\psi & = \frac{\partial \psi}{\partial r} \vect{e}_{r} + \frac{1}{r}\frac{\partial \psi}{\partial \theta} \vect{e}_{\theta} + \frac{1}{r\sin\theta}\frac{\partial \psi}{\partial \phi} \vect{e}_{\phi}                                                                                                            \\
	\Laplasian\psi    & = \frac{1}{r^2} \frac{\partial}{\partial r} \left( r^2 \frac{\partial \psi}{\partial r} \right) + \frac{1}{r^2 \sin \theta} \frac{\partial \psi}{\partial \theta} \left( \sin \theta \frac{\partial \psi}{\partial \theta} \right) + \frac{1}{r^2\sin^2 \theta} \frac{\partial^2 \psi}{\partial \phi^2} \\
	\divg\vect{A}     & =\frac{1}{r^2}\frac{\partial \left(r^2A_r\right)}{\partial r}
	+
	\frac{1}{r\sin \theta }\frac{\partial}{\partial \theta }\left(A_{\theta }\sin \theta \right)
	+
	\frac{1}{r\sin\theta}\frac{\partial A_{\phi}}{\partial \phi }                                                                                                                                                                                                                                                            \\
	\rot\vect{A}      & = \frac{1}{r\sin \theta }\left(\frac{\partial}{\partial \theta }\left(A_{\phi }\sin \theta \right)-\frac{\partial A_{\theta }}{\partial \phi }\right)\vect{e}_{r}
	+
	\frac{1}{r}\left(\frac{1}{\sin \theta }\frac{\partial A_{r}}{\partial \phi }-\frac{\partial}{\partial r}\left(rA_{\phi }\right)\right)\vect{e}_{\theta}
	+ \nonumber                                                                                                                                                                                                                                                                                                                    \\
	                  & + \frac{1}{r}\left(\frac{\partial}{\partial r}\left(rA_{\theta }\right)-\frac{\partial A_{r}}{\partial \theta }\right)\vect{e}_{\phi}
\end{align}
Орти сферичної системи координат зв'язані з декартовими ортами як:
\begin{align}\label{}
    \vect{e}_r &= \sin\theta\cos\phi\, \vect{e}_x + \sin\theta\sin\phi\, \vect{e}_y + \cos\theta\, \vect{e}_z, \\
    \vect{e}_{\theta} &= \cos\theta\cos\phi\, \vect{e}_x + \cos\theta\sin\phi\, \vect{e}_y - \sin\theta\, \vect{e}_z, \\
    \vect{e}_{\phi} &= -\sin\phi\, \vect{e}_x + \cos\phi\, \vect{e}_y. \\
\end{align}

\subsection{Другі похідні}

\begin{align}
	\mathrm{rot}\,\mathrm{grad}\,\phi    & = \rot(\vect{\nabla}\phi)  = 0                                             \\
	\mathrm{div}\,\mathrm{rot}\,\vect{A} & = \divg(\rot\vect{A})  = 0                                                 \\
	\mathrm{rot}\,\mathrm{rot}\,\vect{A} & = \rot(\rot\vect{A})  = \vect{\nabla}(\divg\vect{A}) - \Laplasian \vect{A}
\end{align}


\subsection{Похідні від добутків}

\begin{align}
	\mathrm{grad}\,(\phi \psi)             & = \psi\,\mathrm{grad}\,\phi +\phi\, \mathrm{grad}\,\psi                                                                                                                           \\
	\mathrm{div}\,(\phi \vect{A})          & = \phi\,\mathrm{div}\,\vect{A} + \vect{A}\,\mathrm{grad}\,\phi                                                                                                                    \\
	\mathrm{rot}\,(\phi \vect{A})          & = \phi\,\mathrm{rot}\,\vect{A} + \mathrm{grad}\,\phi \times \vect{A}                                                                                                              \\
	\mathrm{grad}\,(\vect{A}\cdot\vect{B}) & = \vect{B}\times\mathrm{rot}\,\vect{A} + \vect{A}\times\mathrm{rot}\,\vect{B} + \left( \vect{B}\vect{\nabla}\right)\vect{A} + \left( \vect{A}\vect{\nabla}\right)\vect{B}         \\
	\mathrm{div}\,(\vect{A}\times\vect{B}) & = \vect{B}\cdot\mathrm{rot}\,\vect{A} - \vect{A}\cdot\mathrm{rot}\,\vect{B}                                                                                                       \\
	\mathrm{rot}\,(\vect{A}\times\vect{B}) & = \left( \vect{B}\vect{\nabla}\right)\vect{A} - \left( \vect{A}\vect{\nabla}\right)\vect{B} + 	\vect{A}\,\mathrm{div}\,\vect{B} - \vect{B}\,\mathrm{div}\,\vect{A} \label{rotvect} \\
	\frac12\mathrm{grad}\,A^2              & =  \left( \vect{A}\vect{\nabla}\right)\vect{A} + \vect{A}\times\mathrm{rot}\,\vect{A}
\end{align}


\subsection{Інтегральні характеристики та теореми}

Теорема Остроградського-Гаусса:
\begin{equation}\label{OGTheorem}
	\oiint\limits_{\partial \Omega} \vect{A}\cdot d\vect{S} = \iiint\limits_{\Omega} \divg\vect{A} dV,
\end{equation}
де $\Omega$~--- об'єм, $\partial\Omega$ --- його межа.

Теорема Стокса:
\begin{equation}\label{Stoksheorem}
	\oint\limits_{\partial S} \vect{A}\cdot d\vect{l} = \iint\limits_S \rot\vect{A} \cdot d\vect{S},
\end{equation}
де $S$~--- поверхня, що спирається на контур $\partial S$.

Теорема Гріна:
\begin{equation}\label{Grin}
	\iiint\limits_{\Omega}(\phi \nabla ^{2}\psi -\psi \nabla ^{2}\phi )\ dV=\iint \limits_{\partial \Omega}(\phi \vect{\nabla} \psi -\psi \vect{\nabla}
	\phi )\cdot d\vect{S}.
\end{equation}


\section{Поліноми Лежандра}\label{Polinoms}

Поліноми Лежандра застосовуються у теорії потенціалу при розкладанні виразу в околі точки $\vect{r}$:
\begin{equation*}\label{}
    \frac{1}{|\vect{r} - \vect{r}_0|} = \frac{1}{\sqrt{r^2 - 2rr_0 \cos\chi + r_0^2 }}  =  \sum\limits_{l = 0}^{\infty} \frac{r^l_<}{r^{l+1}_>} P_l(\cos\chi),
\end{equation*}
де $r_>$, $r_<$~--- більша і менша із величин $|\vect{r}|$ та $\vect{r}_0$, відповідно, $\cos\chi$~--- кут між векторами $\vect{r}$ та $|\vect{r}_0|$.

\medskip%
\textbf{Деякі поліноми Лежандра}

\begin{align*}
P_{0}(\cos\chi)  = 1, &\quad P_{1}(\cos\chi)  = \cos\chi, \\
P_{2}(\cos\chi)  = \frac {1}{2}(3\cos^2\chi-1), &\quad P_{3}(\cos\chi)  = \frac {1}{2}(5\cos^2\chi-3\cos\chi).
\end{align*}


\section{Сферичні гармоніки}\label{Spherical_Harmonics}

Сферичні функції, що залежать від полярних кутів визначаються формулою:
\begin{equation*}
    Y_{lm}(\theta,\phi) = (-1)^{(m+|m|)/2} \sqrt{\frac{2l+1}{4\pi}\frac{(l-m)!}{(l+m)!}}P_l^{|m|}(\cos \theta) e^{im\phi},
\end{equation*}
де $l = 0,1,2, \ldots$, $m$ пробігає значення від $-l$ до $l$, а $P_l^{|m|}(x)$~--- приєднані функції Лежандра.

Вони утворюють повну ортонормовану систему функцій:
\[
    \int\limits_{\theta=0}^{\pi} \int\limits_{\phi=0}^{2\pi} Y^*_{l',m'}(\theta, \phi) Y_{l,m}(\theta, \phi) \sin\theta d\theta d\phi  = \delta_{l,l'} \delta_{m,m'}.
\]

Деякі сферичні гармоніки:
\begin{align*}
	Y_{0,0}(\theta,\phi)     & =\sqrt{1\over 4\pi},                                            \\
	Y_{1,\pm 1}(\theta,\phi) & =\sqrt{3\over 8\pi} \, \sin\theta \, e^{\pm i\phi},             \\
	Y_{1,0}(\theta,\phi)     & =\sqrt{3\over 4\pi}\, \cos\theta ,                              \\
	Y_{2,0}(\theta,\phi)     & =\sqrt{5\over 16\pi}\, (3\cos^{2}\theta-1),                     \\
	Y_{2,\pm 1}(\theta,\phi) & =\sqrt{15\over 8\pi}\, \sin\theta\, \cos\theta\, e^{\pm i\phi}, \\
	Y_{2,\pm 2}(\theta,\phi) & =\sqrt{15\over 32\pi} \, \sin^{2}\theta \, e^{\pm2i\phi}.
\end{align*}

\section{Циліндричні функції}\setcounter{equation}{0}

Рівняння, що виникають в задачах з циліндричною симетрією, мають вигляд:
	\begin{equation}\label{eq:Bessel_eq}
		\frac{d^2y}{dx^2} + \frac1x\frac{dy}{dx} + \left(1 - \frac{m^2}{x^2} \right) y = 0,
	\end{equation}
	розв'язок яких можна представити за допомогою функцій Бесселя $J_m(x)$ та Неймана $N_m$ у вигляді лінійної комбінації $y(x) = A J_m(x) + B N_m(x)$ або у вигляді лінійної комбінації $y(x) = A H^{(1)}_m(x) + B H^{(2)}_m(x)$, де функції $H_m^{(1,2)} =J_m \pm i N_m$~--- називаються функціями Ганкеля 1-го та 2-го роду, відповідно. Доцільність введення функцій Ганкеля обумовлена тим, що вони мають прості асимптотичні розкладання при $|x| \gg 1$ і зручні для задач, пов'язаних з поширенням хвиль.

Для $m = 0,1,2,\ldots$ функції Неймана нескінченні в точці $x = 0$, тобто $\lim\limits_{x\to0}N_m(x) = -\infty$.

Функції Бесселя можна представити за допомогою ряду (в околі точки $x = 0$ для цілих, або невід'ємних $m$):
\begin{equation}\label{eq:J}
    J_m(x) = \sum_{n=0}^\infty \frac{(-1)^n}{n! \Gamma(n+m+1)} {\left({\frac{x}{2}}\right)}^{2n+m},
\end{equation}
де $ \Gamma$~--- \href{https://en.wikipedia.org/wiki/Gamma_function}{гамма-функція}. Для $m \in \mathbb{Z}$ має місце рівність $J_{-m}(x)  = (-1)^m J_m(x)$.

\begin{center}
  \begin{tikzpicture}
    \begin{axis}[axis lines = middle,
			axis line style={-stealth},
			minor grid style = {line width=.1pt,draw=gray!10},
            width=\textwidth, height=0.5*\textwidth, xlabel=$x$,
            xtick=\empty,
%            ytick={-0.5,0.5,1},
            legend style={draw=none},
    ]
    \addplot+[id=parable,domain=-20:20, samples=500, mark=none, width=2pt, color=red, thick]
    gnuplot{besj0(x)};% node[pin=95:{$J_0(x)$}]{};
    \addplot+[id=parable,domain=-20:20, samples=500, mark=none, width=2pt, color=blue, thick]
    gnuplot{besj1(x)};% node[pin=130:{$J_1(x)$}]{};
    \addplot+[id=parable2,domain=-20:20, samples=500, mark=none, width=2pt, color=green!50!black, thick]
    gnuplot{2*1/x*besj1(x)-besj0(x)};% node[pin=-140:{$J_2(x)$}]{};
    \legend{$J_0(x)$,$J_1(x)$,$J_2(x)$}
   \end{axis}
  \end{tikzpicture}

    {Графіки функцій Бесселя $J_m$ для $m = 0,1,2$.}
\end{center}

\begin{center}
  \begin{tikzpicture}
    \begin{axis}[axis lines = middle,
			axis line style={-stealth},
			minor grid style = {line width=.1pt,draw=gray!10},
            width=\textwidth, height=0.5*\textwidth, xlabel=$x$,
            xtick=\empty,
%            ytick={-0.5,0.5,1},
            legend style={at={(current axis.south east)}, anchor = south east, draw=none},
    ]
    \addplot+[id=parable,domain=0:40, restrict y to domain=-1.5:1, samples=500, mark=none, width=2pt, color=red, thick]
    gnuplot{besy0(x)};% node[pin=95:{$J_0(x)$}]{};
    \addplot+[id=parable,domain=0:40, restrict y to domain=-1.5:1, samples=1000, mark=none, width=2pt, color=blue, thick]
    gnuplot{besy1(x)};% node[pin=130:{$J_1(x)$}]{};
    \addplot+[id=parable2,domain=0:40, restrict y to domain=-1.5:1, samples=1000, mark=none, width=2pt, color=green!50!black, thick]
    gnuplot{2*1/x*besy1(x)-besy0(x)};% node[pin=-140:{$J_2(x)$}]{};
    \legend{$N_0(x)$,$N_1(x)$,$N_2(x)$}
   \end{axis}
  \end{tikzpicture}

    {Графіки функцій Неймана $N_m$ для $m = 0,1,2$.}
\end{center}

Функції Неймана визначаються через функції Бесселя як:
\begin{equation}\label{eq:NG}
    N_m(x) = \frac{J_m(x)\cos m\pi - J_{-m}(x)}{\sin m\pi}.
\end{equation}


Деякі рекурентні співвідношення:
\begin{equation}
    J_{m + 1}(x) + J_{m - 1}(x) = \frac{2m}{x}J_m(x).
\end{equation}

Деякі диференціальні та інтегральні співвідношення для нецілих $m$ (для цілих $m$ ці функції можна визначити за допомогою граничного переходу).:
\begin{align}
    \frac{d}{dx} J_0(x) = - J_1(x), \\
    \frac{d}{dx} \left( x^{-m}J_m(x)\right)  = - x^{-m}J_{m+1}(x), \\
     \int\limits_0^{x} x^{\prime m+1} J_m(x') dx' &= x^{m+1}J_{m+1}. \label{eq:recInt}
\end{align}

Інтеграли від добутків:
\begin{equation}
    \int\limits_0^x  J_m(k_1x')J_m(k_2x') x' dx' = \frac{x\left( k_2J_m(k_1x)J'_m(k_2x) - k_1J_m(k_2x)J'_m(k_1x)\right)}{k_1^2-k_2^{2}}   \label{eq:JJ0*}.
\end{equation}

В задачах, зазвичай, часто необхідно знайти наближений вигляд циліндричних функцій при малих та великих значеннях аргументу $x$:

при $|x| \ll 1$ з~\eqref{eq:J}

\begin{equation}
    J_0(x) \approx 1 - \frac{x^2}{4}, \quad
    J_m \approx \frac{x^m}{2^m m!}, \ m \ge 1,\, m \in \mathbb{N};
\end{equation}

при $|x| \gg 1$

\begin{align}
    J_m &\approx \sqrt{\frac{2}{\pi x}}  \cos\left( x - m\frac{\pi}{2} - \frac{\pi}{4}\right),
      \label{eq:Jxgg1}\\
    N_m &\approx \sqrt{\frac{2}{\pi x}}  \sin\left( x - m\frac{\pi}{2} - \frac{\pi}{4}\right),
  \label{eq:Yxgg1}\\
    H_m^{(1,2)} &\approx \sqrt{\frac{2}{\pi x}}  e^{\pm i \left( x - m\frac{\pi}{2} - \frac{\pi}{4}\right) }. \label{eq:Hxgg1}
\end{align}

Співвідношення Якобі-Ангера (розкладання за функціями Бесселя):
\begin{equation}\label{eq:JacobiAnger}
    e^{ix\cos\theta} = \sum\limits_{m= - \infty}^{\infty} i^mJ_m(x)e^{im\theta}, \quad
    e^{ix\sin\theta} = \sum\limits_{m= - \infty}^{\infty} J_m(x)e^{im\theta}. \\
\end{equation}

\section{Дельта-функція Дірака}\setcounter{equation}{0}

Дельта-функція Дірака (або $\delta$-функція) є узагальненою функцією і була введена фізиком Полем Діраком для моделювання густини ідеалізованої точкової маси або точкового заряду.

На <<фізичному рівні строгості>> можна визначити $\delta$-функцію формальним співвідношенням:
\begin{equation}
    \int\limits_{-\infty}^{+\infty}f(x)\delta(x-x_0)\,dx=f(x_0),
\end{equation}
У випадку інтегрування по скінченному об'єму $V$:
\begin{equation}\label{eq:delta3D}
    \int\limits_{V}f(\vect{r})\delta(\vect{r} - \vect{r}_0) dV=f(\vect{r}_0),
\end{equation}
де точка $\vect{r}_0$ знаходиться всередині об'єму $V$.

%$\delta$-функцію однієї дійсної змінної можна визначити як функцію, що задовольняє наступним умовам:
%\begin{equation}\label{eq:delta}
%\delta(x - x_0)=\left\{\begin{matrix}
%   +\infty, & x=x_0, \\
%   0, & x \neq x_0; \\
%\end{matrix}\right.
%\end{equation}
%\begin{equation}\label{eq:deltaProp0}
%    \int\limits_{-\infty}^{+\infty}\delta(x-x_0)dx=1.
%\end{equation}
%Тобто ця функція не дорівнює нулю тільки в точці $x = x_0$, де вона перетворюється в нескінченність таким чином, щоб її інтеграл в будь-якому околі точки $x = x_0$ дорівнює $1$.
%
%Аналітичне представлення $\delta$-функції:
%\begin{equation}\label{eq:deltaanalit}
%    \delta(x) = \frac{1}{2\pi}\int\limits_{-\infty}^{\infty} e^{ikx}dk.
%\end{equation}
\bigskip\noindent%
\textbf{Властивості дельта-функції}
\bigskip

\begin{enumerate}[label=\alph*)]
\item Дельта-функція парна $\delta(-x) = \delta(x)$,
\item $x\delta(x) = 0$,
\item $\delta(ax) = \frac{1}{|a|}\delta(x)$,
\item  $x\delta^\prime(x)=-\delta(x)$,
\item $\delta(f(x))=\sum\limits_k\frac{\delta(x-x_k)}{|f'(x_k)|}$, де $x_k$~--- нулі функції $f(x)$,

\end{enumerate}

%\subsection{Розкладання дельта-функції в ряд Фур'є}
%
%$\delta$-Функцію можна представити у вигляді ряду Фур'є на деякому проміжку $x \in (-l,l)$:
%\begin{equation}\label{eq:deltacos}
%    \delta(x) = \frac{1}{2l} +  \frac1l \sum\limits_{m=1}^{\infty}\cos\frac{m\pi}{l}x,
%\end{equation}
%або в комплексній формі
%\begin{equation}\label{eq:deltae}
%    \delta(x) =  \frac{1}{2l} \sum\limits_{m=-\infty}^{\infty}e^{i\frac{2m\pi}{l}x}.
%\end{equation}
%У випадку, якщо ми маємо періодичну функцію у вигляді $\delta$-<<частоколу>>~--- $\sum\limits_{k=-\infty}^{\infty}\delta(x-ka)$~--- відстань між піками якого дорівнює $a=2l$, то для нього розкладання також матиме вигляд~\eqref{eq:deltacos} (або \eqref{eq:deltae}):
%\begin{equation}\label{eq:deltalll}
%    \sum\limits_{k=-\infty}^{\infty}\delta(x-ka)  = \frac{1}{a} +  \frac2a \sum\limits_{m=1}^{\infty}\cos\frac{2m\pi}{a}x
%\end{equation}
% і буде справедливе на всій числовій прямій $x \in (-\infty,\infty)$.
%
%\subsection{Тривимірна дельта-функція}
%
%Тривимірна $\delta$-функція визначається співвідношеннями:
%\begin{equation}\label{eq:delta3D}
%    \delta(\vect{r} - \vect{r}_0) = \delta(x-x_0)\delta(y-y_0)\delta(z-z_0),
%\end{equation}
%\begin{equation}\label{eq:delta3D}
%    \int\limits_{V}f(\vect{r})\delta(\vect{r} - \vect{r}_0) dV=f(\vect{r}_0),
%\end{equation}
%де точка $\vect{r}_0$ знаходиться в середині об'єму $V$.
%
%Аналітичне представлення тривимірної $\delta$-функції:
%\begin{equation}\label{eq:deltaanalit3В}
%    \delta(\vect{r}) = \frac{1}{(2\pi)^3}\int\limits_{-\infty}^{\infty} e^{i\vect{k}\cdot\vect{r}}d\vect{k},
%\end{equation}
%де $d\vect{k}=dk_xdk_ydk_z$.



%% !TeX program = lualatex
% !TeX encoding = utf8
% !TeX spellcheck = uk_UA
% !TeX root =../FTProblems.tex

\newpage
\section{Фізичні константи}

\begin{table}[h!]\centering
	\tabcaption{Деякі фізичні константи в системі СГС}
	\label{tab:SItoGauss}
	\small
	\begin{tabular}{lcl}
		\toprule
		Константа                  & Символ   & \centercell{Значення}                           \\ \midrule
		Швидкість світла у вакуумі & $c$      & $2.997 924 58 \cdot 10^{10}$ см/с               \\
		Гравітаційна стала         & $G$      & $6.674 28 \cdot 10^{-8}$ см$^3$/(г$\cdot$c$^2$) \\
		Стала Планка               & $\hbar$  & $1.054 5716 \cdot 10^{-27}$ ерг$\cdot$с         \\
		Елементарний заряд         & $e$      & $4.803 204 27  \cdot 10^{-10}$ СГСЕ$_q $\footnotemark        \\
                                   &          & $1.60217662 \cdot 10^{-19}$ Кл (в SI)           \\ 
		Маса електрона             & $m_e$    & $9.109 382 15 \cdot 10^{-20}$ г                 \\
		Енергія спокою електрона   & $m_ec^2$ & $0.511 \cdot 10^{6}$ еВ                         \\
		Маса протона               & $m_p$    & $1.672 621 9 \cdot 10^{-30}$ г                  \\
		Енергія спокою протона     & $m_pc^2$ & $938.26 \cdot 10^{6}$ еВ                        \\
		\hskip4ex Електрон-Вольт   & еВ       & $1.602\cdot 10^{-12}$ ерг                       \\
		Борівський радіус          & $a_0$    & $5.291 772 0859 \cdot 10^{-9}$ см               \\
		Магнетон Бора              & $\mu_B$  & $9.274 009 15 \cdot 10^{-21}$ ерг/Гс            \\
		Стала Больцмана            & $k$      & $1.380 6504 \cdot 10^{-16}$ ерг/К               \\
		Універсальна газова стала  & $R$      & $8.314\cdot 10^{7}$ ерг/(К$\cdot$моль)          \\
		Число Авогадро             & $N_A$    & $6.022\cdot 10^{23}$  моль$^{-1}$               \\
		Атомна одиниця маси        & $u$      & $1.66042 10^{-24}$  г                           \\ \bottomrule
	\end{tabular}
\end{table}

\footnotetext{Міжнародна назва цієї одиниці~--- Франклін (позначення: Fr, Фр). У радянській і сучасній вітчизняній літературі зазвичай називається просто СГСЕ-од. заряду. Франклін є однією з чотирьох основних одиниць системи Гауса (поряд з сантиметром, грамом і секундою).}

%%========================================================================================================
%
%									       Вставка бібліографії
%
%=========================================================================================================
\clearpage\pagestyle{bibliography}

\nocite{%
	ZhdanovFT,
	ZhdanovESS,
    ZhdanovRelativity,
	TerletskyElectroDyn,
	Jackson,
	GriffithsElectro,
	PennerUgarov,
	Tamm,
	GreinerElectrodynamics,
	LL2,
	LL8,
	Vlasov,
	% ==========================
	BatyginElectroRus,
	Grechko,
	Kramm,
	Zhyrnov,
	AlekseevElectro,
	VekshteinElectro,
	Ginldenburg,
    ZangwillElectro,
}



\printbibheading[heading=bibintoc]

\defbibheading{subheading}[\bibname]{%
	\section*{#1\label{BibPage}}
}

\printbibliography[category=Textbooks, heading=subheading, title={Підручники та посібники}]
\addtocategory{Textbooks}{%
	ZhdanovFT,
	ZhdanovESS,
    ZhdanovRelativity,
	Jackson,
	GriffithsElectro,
	PennerUgarov,
	Tamm,
	GreinerElectrodynamics,
	LL2,
	TerletskyElectroDyn,
	Vlasov,
	LL8,
    ZangwillElectro,
}


\printbibliography[category=Problems, heading=subheading, title={Задачники}]
\addtocategory{Problems}{%
	BatyginElectroRus,
	Grechko,
	Kramm,
	Zhyrnov,
	AlekseevElectro,
	VekshteinElectro,
	Ginldenburg,
}

%\makelastpage
\end{document}
